\documentclass[10pt,a4paper]{amsart}
\usepackage[margin=19mm]{geometry}
\usepackage{amsmath,amssymb,mathtools}
\usepackage[scr=rsfs,cal=euler]{mathalfa}
\usepackage{xfrac}
\usepackage[unicode,hidelinks,bookmarks=false]{hyperref}
\newcommand{\ie}{i.e.\ }
\newcommand{\cf}{cf.\ }
\newcommand{\resp}{respectively\ }
\newcommand{\Subsection}{\subsection}
\providecommand{\nobreakdash}{\nobreak\hbox{-}}
\setlength{\emergencystretch}{2em}
\multlinegap0pt
\usepackage{bm}
\usepackage{bbm}
\usepackage{dsfont}
\usepackage{xspace}
\usepackage{cancel}
\usepackage{mathtools}
\usepackage{amsthm}
\makeatletter
\@ifundefined{theorem}{\newtheorem{theorem}{Theorem}}{}
\@ifundefined{definition}{\newtheorem{definition}[theorem]{Definition}}{}
\@ifundefined{prop}{\newtheorem{prop}[theorem]{Proposition}}{}
\makeatother
\DeclareMathOperator{\R}{\mathbb{R}}
\newcommand{\LH}{\operatorname{LH}_\infty^1}
\newcommand{\Lp}{L^{\p}(\R)}
\newcommand{\defeq}{\vcentcolon=}
\newcommand{\p}{p(\cdot)}
\title[The Fourier transform in variable exponent Lebesgue spaces]{The Fourier transform in variable exponent Lebesgue spaces}
\author{Andr\'e Pedroso Kowacs, Wagner Augusto Almeida de Moraes}
\date{Source: arXiv:2506.08891v1 (CC BY 4.0)}
\begin{document}
\maketitle

\noindent\emph{Source and licence.} The Fourier transform in variable exponent Lebesgue spaces. Version arXiv:2506.08891v1 (CC BY 4.0), distributed under the Creative Commons Attribution 4.0 International licence, as recorded in the arXiv licence metadata for this exact version and shown on the abstract page https://arxiv.org/abs/2506.08891v1. Full rights notice: licenses/arxiv-2506.08891-CC-BY-4.0.txt.\par\smallskip

\noindent\textit{Editorial scope.} Counted unit: the Prop whose source label is \texttt{prop\_properties} in the article (source0.tex lines 1084--1096, source label \texttt{prop\_properties}), delivered together with the article's own complete original proof of it (source0.tex lines 1097--1114, one \texttt{proof} environment of 18 lines). No mathematics is written, repaired or re-derived: statement and proof are the article's own text line for line. Required context retained uncounted: theo\_1 (lines 935--942); definition\_Fourier (lines 991--995). Every definition, display and label of those lines is retained unchanged. Omitted from this package and not redistributed: 1--934, 943--990, 996--1083, 1115--1390. Nothing inside a retained excerpt is omitted or altered: every retained body is delivered exactly as the article's lines are written, so no omission receipt of any kind accompanies this package. Citations: the retained excerpts carry no citation command at all, so no bibliography entry of the article is transcribed and no reference is redistributed. Reference adaptation: none was needed and none was made; every reference inside the delivered unit resolves inside it, so no unresolved reference or citation is produced. Printed slips kept literal: no printed word, symbol, hypothesis or number of the counted statement or of its proof was altered. Layout and class: the article's own class is \texttt{amsart} with packages including inputenc, amsmath, amssymb, wrapfig, url, mathtools, graphicx, stmaryrd, amsthm, xcolor, hyperref, enumitem, comment, relsize; the wrapper is the lane's pinned \texttt{amsart} editorial wrapper at 10pt on A4, which restates the theorem-like environments the retained text needs and adds the article's own macro definitions that the retained text needs (\texttt{\textbackslash LH}, \texttt{\textbackslash Lp}, \texttt{\textbackslash R}, \texttt{\textbackslash defeq}, \texttt{\textbackslash p}), with extra wrapper packages (bm, bbm, dsfont, xspace, cancel, mathtools). The delivered numbering is the unit's own. Assets: no retained span contains a figure environment, a graphics-inclusion command, a PDF-page inclusion, a table or a TikZ picture; no image file is copied into this package, the package declares no asset dependency and no image file is redistributed with it. Licence: version arXiv:2506.08891v1 of this article is distributed under the Creative Commons Attribution 4.0 International licence, as recorded in the arXiv licence metadata for this exact version and shown on the abstract page https://arxiv.org/abs/2506.08891v1. A full rights notice is delivered at licenses/arxiv-2506.08891-CC-BY-4.0.txt. Characters: every retained excerpt is pure ASCII, so the delivered engine, XeLaTeX with its default Latin Modern text font, reports no missing character.
\begin{theorem}\label{theo_1}
     Let $\p\in \mathcal{P}(\R)$ satisfy $p_{+}<\infty$ or $1/\p$ is $\LH$.
     \begin{enumerate}
         \item Define $\mathcal{B}_{p(\cdot)}(\R) = \{\Psi_f : f\in L^{p(\cdot)}(\R)\}$. For $f\in L^{\p}(\R)$, let $\|\Psi_f\|_{\mathcal{B}_{\p}(\R)}\defeq\|f\|_{\p}$. Then  $\left(\mathcal{B}_{p(\cdot)}(\R),\|\cdot\|_{\mathcal{B}_{\p}(\R)}\right)$ is a Banach space isometrically isomorphic to $L^{\p}(\R)$.
         \item Define $\mathcal{A}_{\p}(\R) = \{\Psi_f': f\in L^{\p}(\R)\}$, where the superscript $'$ denotes the distributional derivative. For $f\in L^{\p}(\R)$, let $\|\Psi_f'\|_{\mathcal{A}_{\p}(\R)}\defeq\|f\|_{\p}$. Then 
         \noindent $\left(\mathcal{A}_{\p}(\R),\|\cdot\|_{\mathcal{A}_{\p}(\R)}\right)$ is a Banach space isometrically isomorphic to $\Lp$.
     \end{enumerate}
\end{theorem}

\begin{definition}\label{definition_Fourier}
    Consider $\p\in\mathcal{P}(\R)$ such that $p_{+}<\infty$ or $1/\p$ is $\LH$. Define the Fourier transform as $\mathcal{F}:\Lp\to \mathcal{A}_{\p}(\R)$ by $\mathcal{F}(f)\equiv\widehat{f}\defeq\Psi_f'$.
    Or equivalently,
    $$\langle \widehat{f},\varphi\rangle=-\langle \Psi_f,\varphi'\rangle=-\int_{-\infty}^{\infty}\Psi_f(s)\varphi'(s)ds,$$ for each $\varphi\in\mathcal{S}(\R)$. The space of Fourier transforms on $L^{\p}(\R)$ is then denoted by $\mathcal{A}_{\p}(\R)$. 
\end{definition}

\begin{prop}\label{prop_properties}
    Consider $\p\in\mathcal{P}(\R)$ such that $p_{+}<\infty$ or $1/\p$ is $\LH$ and fix  $f\in\Lp$. For any function $g:\R\to\mathbb{C}$, denote $\tau_ag(x)=g(x-a)$ and $\tilde{g}(x)=g(-x)$, for every $a,x\in\R$. Then
    \begin{enumerate}
               \item  The Fourier transform is a bounded linear isomorphism and isometry $\mathcal{F}$ : $L^{\p}(\mathbb{R}) \to \mathcal{A}_{\p}(\mathbb{R})$, that is, $\|\hat{f}\|_{\mathcal{A}_{\p}(\R)}=\|f\|_{\p}$;

\item $\hat{f}^{(n)}=\Psi_f^{(n+1)}$;
\item For $a \in \mathbb{R}$, we have that  $\widehat{\tau_a f}=\Psi_{\tau_a f}'$, and $\Psi_{\tau_a f}(s)=f * \tilde{u}_s(-a)$;
\item For $a \in \mathbb{R}$ and $F(t)={e}^{{i} a t} f(t)$, we have that $\widehat{F}=\left(\tau_a \Psi_f\right)^{\prime}$;
\item $\hat{\tilde{f}}=\Psi_{\tilde{f}}^{\prime}$, and $\Psi_{\tilde{f}}=-\widetilde{\Psi}_f$;
\item If $\hat{f} \in L^{r(\cdot)}(\mathbb{R})$ for some $r(\cdot)\in \mathcal{P}(\R)$ such that $r_{+}<\infty$ or $1/r(\cdot)$ is $\LH$, then $\hat{\hat{f}}=\Psi_{\widehat{f}}'$;
\item Let $g(x)=f(a x+b)$ for $a, b \in \mathbb{R}$ with $a \neq 0$. Then $\hat{g}=\Psi_g^{\prime}$, where $\Psi_g(s)=$ $\operatorname{sgn}(a) \Psi_{\tau_{-b} f}(s / a)$.
    \end{enumerate}
\end{prop}

\begin{proof}
    Claim {\it (1)} follows from Definition \ref{definition_Fourier} and Theorem \ref{theo_1}.
Claims \noindent{\it (2), (3)} and {\it (6)} 
%$\left\langle D^n \hat{f}, \varphi\right\rangle=\left\langle\Psi_f^{(n+1)}, \varphi\right\rangle=(-1)^{n+1}\left\langle\Psi_f, \varphi^{(n+1)}\right\rangle$. Integration by parts gives another formula for the derivative of a Fourier transform.
follow immediately from the definitions.

\noindent For the proof of {\it (4)}, notice that
$$
\Psi_F(s)=\int_{-\infty}^{\infty}\left(\frac{1-{e}^{-{i}(s-a) t}}{{i} t}-\frac{1-{e}^{{i} a t}}{{i} t}\right) f(t) {d} t=\Psi_f(s-a)-\Psi_f(-a),
$$
therefore $\widehat{F}=\left(\tau_a \Psi_f\right)^{\prime}$.

\noindent For {\it (5)}, note that from the definition of $\tilde{f}$ and performing a change of variables, we have that
$$
\Psi_{\tilde{f}}(s)=\int_{-\infty}^{\infty}\left(\frac{1-{e}^{-{i} s t}}{{i} t}\right) f(-t) {d} t=-\int_{-\infty}^{\infty}\left(\frac{1-{e}^{{i} s t}}{{i} t}\right) f(t) {d} t.
$$
\noindent Finally, {\it (7)} follows from a change of variables in the integral which defines $\Psi_g$.
\end{proof}

\end{document}
